\section*{Zusammenfassung} 
Diese Studienarbeit befasst sich mit der Weiterentwicklung eines automatisierten
Modell-Hochregallagers an der Dualen Hochschule Baden-Württemberg. Ziel ist es,
bestehende Funktionen um sichere und strukturierte Erweiterungen zu ergänzen und
das System stärker an industrielle Anforderungen anzupassen. Der Fokus liegt auf der
Integration einer sicherheitsgerichteten Lichtschranke, der Umsetzung eines internen
Umlagerungsprozesses sowie der Optimierung der SPS-Programmstruktur durch den
Einsatz der Programmiersprache Structured Text (ST). Die Lichtschranke wird über
IO-Link an eine Siemens S7-1500 angebunden und mithilfe geeigneter Prozesswerte
in die Sicherheitslogik integriert. Durch eine gezielte Auswertung des Prozesswertes
„Number of Beams Occupied“ kann der Gefahrenbereich zuverlässig überwacht und bei
unzulässigen Eingriffen eine automatische Abschaltung ausgelöst werden. Die Sicher
heitsfunktion ist vollständig in die bestehende Steuerungsarchitektur eingebettet und
mit dem Störungsmanagement gekoppelt. Ergänzend wird ein Umlagerungsprozess
realisiert, der die interne Verlagerung von Lagergütern ermöglicht. Eine eigenständige
Schrittkette stellt einen sicheren und flexiblen Ablauf des Prozesses sicher und berück
sichtigt sowohl die Bewegungslogik der Achsen als auch das Lagerplatzmanagement.
Insgesamt wird das Hochregallagersystem funktional, sicherheitstechnisch und struk
turell erweitert und bildet eine belastbare Grundlage für zukünftige, auftragsbasierte
Erweiterungen